\documentclass[11pt,a4paper]{report}
\usepackage[utf8]{inputenc}
\usepackage[T1]{fontenc}
\usepackage[french,english]{babel}
\usepackage{amsmath,amssymb,amsfonts}
\usepackage{graphicx}
\usepackage{booktabs}
\usepackage{geometry}
\usepackage{hyperref}
\usepackage{xcolor}
\usepackage{float}
\usepackage{listings}
\usepackage{tikz}
\usepackage{caption}
\usepackage{subcaption}

\geometry{margin=2.5cm}
\hypersetup{
    colorlinks=true,
    linkcolor=blue!70!black,
    citecolor=blue!70!black,
    urlcolor=teal
}

\definecolor{codegray}{rgb}{0.5,0.5,0.5}
\definecolor{codepurple}{rgb}{0.58,0,0.82}
\definecolor{backcolour}{rgb}{0.96,0.96,0.96}

\lstdefinestyle{mystyle}{
    backgroundcolor=\color{backcolour},
    commentstyle=\color{codegray},
    keywordstyle=\color{blue},
    numberstyle=\tiny\color{codegray},
    stringstyle=\color{codepurple},
    basicstyle=\ttfamily\footnotesize,
    breaklines=true,
    numbers=left,
    numbersep=5pt
}
\lstset{style=mystyle}

\begin{document}

% =========================================================================
% TITLE PAGE
% =========================================================================
\begin{titlepage}
\centering
{\scshape\LARGE République Tunisienne\par}
{\scshape\large Ministère de l'Enseignement Supérieur et de la Recherche Scientifique\par}
\vspace{1.5cm}

{\scshape\Large Module : Deep Learning Avancé \& MLOps\par}
{\large Cycle Ingénieur en Intelligence Artificielle \& Data Science --- 3\textsuperscript{ème} Année\par}
\vspace{1.5cm}

\rule{\linewidth}{0.5mm}\\[0.4cm]
{\huge\bfseries ArchiGen 3D : Plateforme Industrielle de Synthèse Architecturale et d'Extrusion 3D par Vision Profonde et cGAN\par}
\rule{\linewidth}{0.5mm}\\[1.5cm]

\begin{minipage}{0.48\textwidth}
\begin{flushleft} \large
\emph{Membres de l'Équipe :}\\
\textbf{Khalil Ayeb}\\
\textbf{Amine Ben Fadhel}\\
\textbf{Khawla Hawari}
\end{flushleft}
\end{minipage}
\hfill
\begin{minipage}{0.48\textwidth}
\begin{flushright} \large
\emph{Enseignant Responsable :}\\
\textbf{Prof. Haythem Ghazouani}
\end{flushright}
\end{minipage}

\vfill
{\large Année Universitaire 2026--2027\par}
\end{titlepage}

% =========================================================================
% ABSTRACT / RÉSUMÉ
% =========================================================================
\selectlanguage{french}
\begin{abstract}
Ce projet présente \textbf{ArchiGen 3D}, un système complet de Computer Vision et d'Intelligence Artificielle Générative dédié à la conversion automatisée de plans 2D d'architectes en environnements 3D interactifs avec textures de façades synthétisées. Le système s'articule autour de deux axes fondamentaux : premièrement, une architecture multi-tâches UNet++ (ResNet-34) entraînée sur CubiCasa5K (4\,200 plans d'entraînement, 400 de test, licence CC BY-NC 4.0) assurant la segmentation conjointe des cloisons structurelles (murs intérieurs et extérieurs, portes, fenêtres) et des pièces fonctionnelles (11 classes sémantiques), suivie d'une vectorisation géométrique rectilinéaire et d'une extrusion 3D étanche. Deuxièmement, un modèle génératif antagoniste conditionnel (cGAN / Pix2Pix) stabilisé par normalisation spectrale, perte perceptuelle LPIPS (SqueezeNet) et matching de caractéristiques multi-échelles, synthétisant des textures de façades sur le jeu de données CMP Facades (400 paires train, 106 test). L'intégralité du cycle de vie respecte les standards MLOps les plus stricts : versioning des données et poids par DVC, observabilité complète via MLflow, moteur d'inférence accéléré sous ONNX Runtime avec parité numérique stricte ($\Delta_{\max} = 6.31 \times 10^{-6}$), conteneurisation Docker Compose multi-services avec proxy inverse Nginx, et pipeline CI/CD automatisé sous GitHub Actions validé par 18 tests Pytest.
\end{abstract}

\tableofcontents
\newpage

% =========================================================================
% CHAPITRE 1: INTRODUCTION & CADRAGE INDUSTRIEL
% =========================================================================
\chapter{Introduction \& Contexte Industriel}

\section{Mise en Contexte et Problématique Métier}
Dans l'industrie contemporaine de l'architecture, de l'immobilier et de la modélisation des données du bâtiment (BIM - \emph{Building Information Modeling}), la conversion d'esquisses ou de plans 2D scannés en maquettes numériques 3D exploitables représente une charge de travail chronophage et coûteuse, majoritairement réalisée à la main par des dessinateurs-projeteurs. Bien que des solutions commerciales existent (telles que Matterport ou HomeByMe), la plupart reposent sur des capteurs LiDAR embarqués ou sur des interfaces de saisie manuelles.

L'objectif du projet \textbf{ArchiGen 3D} est de concevoir et d'industrialiser une chaîne de traitement logicielle entièrement automatisée (\emph{Plan-to-3D-House}), recevant en entrée une simple image matricielle 2D (PNG, JPEG) d'un plan d'étage, et produisant en sortie :
\begin{enumerate}
    \item Une cartographie sémantique complète identifiant la nature des espaces (chambre, séjour, cuisine, salle d'eau, etc.) et les éléments de franchissement (portes et fenêtres).
    \item Une géométrie vectorielle vectorisée en coordonnées métriques réelles avec calibration d'échelle.
    \item Un maillage volumique 3D étanche prêt pour le Web (Three.js WebGL) et exportable au format universel OBJ/MTL.
    \item La synthèse générative de textures de façades extérieures architecturales conditionnées par la volumétrie et les ouvertures détectées.
\end{enumerate}

\section{Données et Architectures des Modèles}
Le projet repose sur la complémentarité de deux modèles d'apprentissage profond distincts :
\begin{itemize}
    \item \textbf{Modèle 1 --- Segmentation Multi-Tâches de Plans (MultiTask UNet++) :}
    \begin{itemize}
        \item \emph{Jeu de Données :} CubiCasa5K ($4\,200$ plans d'entraînement, $400$ de validation, $400$ de test sous licence académique CC BY-NC 4.0).
        \item \emph{Architecture :} Réseau UNet++ avec encodeur ResNet-34 pré-entraîné ImageNet et deux têtes de décodage convolutives parallèles : (1) Tête structurelle 5 classes (fond, mur intérieur, mur extérieur, porte, fenêtre) et (2) Tête sémantique 12 classes (11 types de pièces fonctionnelles + fond).
        \item \emph{Objectif :} Minimisation conjointe d'une perte hybride Focal Loss + Dice Loss pour contrer le fort déséquilibre de classe des ouvertures.
    \end{itemize}
    \item \textbf{Modèle 2 --- Synthèse de Textures de Façades (Pix2Pix cGAN Perceptuel) :}
    \begin{itemize}
        \item \emph{Jeu de Données :} CMP Facades ($400$ paires d'entraînement, $106$ paires de test).
        \item \emph{Architecture :} Générateur UNet-256 avec sur-échantillonnage par interpolation bilinéaire et convolution (supprimant les artefacts de damier), Discriminateur PatchGAN $70 \times 70$ stabilisé par normalisation spectrale.
        \item \emph{Objectif :} Perte composite combinant Moindres Carrés (LSGAN), distance $\mathcal{L}_1$ pondérée, perte de correspondance de caractéristiques (\emph{Feature Matching}) et distance perceptuelle LPIPS (SqueezeNet).
    \end{itemize}
\end{itemize}

\newpage

% =========================================================================
% CHAPITRE 2: FORMULATION MATHÉMATIQUE & ARCHITECTURES
% =========================================================================
\chapter{Formulation Mathématique \& Architectures}

\section{Axe Majeur : Synthèse Générative de Façades (cGAN)}
L'axe principal retenu pour l'évaluation officielle est l'\textbf{Axe 2 (Intelligence Artificielle Générative en Imagerie)}. La tâche consiste en une traduction d'image à image conditionnelle (\emph{pix2pix}), apprenant la distribution conditionnelle $p(y|x)$ où $x \in \mathbb{R}^{H \times W \times 3}$ représente la carte sémantique d'élévation et $y \in \mathbb{R}^{H \times W \times 3}$ la façade photographique réelle.

\subsection{Objectif Adversarial : LSGAN vs Minimax Standard}
Pour stabiliser l'apprentissage et éviter l'évanouissement des gradients propre à la perte minimax saturante basée sur la divergence de Jensen-Shannon, nous adoptons la formulation des Moindres Carrés (LSGAN - \emph{Least Squares GAN}, Mao et al., 2017) :
\begin{equation}
    \mathcal{L}_{\text{LSGAN}}(D) = \frac{1}{2} \mathbb{E}_{(x, y)} \left[ (D(x, y) - 1)^2 \right] + \frac{1}{2} \mathbb{E}_{x} \left[ (D(x, G(x)))^2 \right]
\end{equation}
\begin{equation}
    \mathcal{L}_{\text{LSGAN}}(G) = \frac{1}{2} \mathbb{E}_{x} \left[ (D(x, G(x)) - 1)^2 \right]
\end{equation}

\subsection{Fonction de Perte Composée Multi-Critères}
Le générateur $G$ est entraîné à minimiser une fonction de coût hybride combinant réalisme adversarial, cohérence spatiale pixel-à-pixel, matching structurel multi-échelles et fidélité perceptuelle :
\begin{equation}
    \mathcal{L}_{\text{total}}(G) = \mathcal{L}_{\text{LSGAN}}(G) + \lambda_{\mathcal{L}_1} \mathcal{L}_1(G) + \lambda_{\text{FM}} \mathcal{L}_{\text{FM}}(G, D) + \lambda_{\text{perc}} \mathcal{L}_{\text{perc}}(G)
\end{equation}

Les composantes s'établissent comme suit :
\begin{enumerate}
    \item \textbf{Erreur absolue pixel-à-pixel :}
    \begin{equation}
        \mathcal{L}_1(G) = \mathbb{E}_{(x, y)} \left[ \| y - G(x) \|_1 \right]
    \end{equation}
    \item \textbf{Perte d'appariement de caractéristiques (\emph{Feature Matching Loss}) :}
    Inspirée de Pix2PixHD (Wang et al., 2018), cette perte contraint le générateur à produire des activations identiques à celles des images réelles au sein des $T$ couches intermédiaires du discriminateur $D$ :
    \begin{equation}
        \mathcal{L}_{\text{FM}}(G, D) = \mathbb{E}_{(x, y)} \left[ \sum_{i=1}^{T} \frac{1}{N_i} \| D^{(i)}(x, y) - D^{(i)}(x, G(x)) \|_1 \right]
    \end{equation}
    \item \textbf{Perte perceptuelle LPIPS (Learned Perceptual Image Patch Similarity) :}
    Calculée sur un réseau SqueezeNet pré-entraîné sur ImageNet, pénalisant les moyennes floues et imposant des contrastes haute fréquence réalistes :
    \begin{equation}
        \mathcal{L}_{\text{perc}}(G) = \mathbb{E}_{(x, y)} \left[ \sum_{l} \frac{1}{H_l W_l} \| \phi_l(y) - \phi_l(G(x)) \|_2^2 \right]
    \end{equation}
\end{enumerate}

Dans notre configuration de production, les hyperparamètres réels sont rigoureusement fixés à : $\lambda_{\mathcal{L}_1} = 30.0$, $\lambda_{\text{FM}} = 10.0$, $\lambda_{\text{perc}} = 10.0$.

\subsection{Architectures des Réseaux et Éradication des Artefacts}
\begin{itemize}
    \item \textbf{Générateur UNet Anti-Checkerboard :} Contrairement aux couches \texttt{ConvTranspose2d} classiques sujettes aux artefacts de grille d'Odena (\emph{checkerboard artifacts}), notre décodeur emploie un suréchantillonnage par plus proche voisin (\texttt{Nearest-Neighbor Upsample}) suivi d'une convolution 2D classique.
    \item \textbf{Discriminateur PatchGAN avec Normalisation Spectrale :} Opérant sur des patchs réceptifs de $70 \times 70$ pixels, les poids de chaque couche convolutive sont contraints par Spectral Normalization ($\sigma(W) \le 1$), garantissant la régularité 1-Lipschitzienne et prévenant l'effondrement de mode.
\end{itemize}

\section{Axe Complémentaire : Segmentation Multi-Tâches de Plans (UNet++)}
Pour extraire la topologie intérieure, nous mettons en œuvre une architecture UNet++ multi-tâches reposant sur un encodeur convolutif ResNet-34 pré-entraîné, couplé à deux têtes de décodage spécialisées :
\begin{itemize}
    \item \textbf{Tête Structurelle (5 classes) :} Fond, mur intérieur (\texttt{Wall\_Internal}), mur extérieur (\texttt{Wall\_External}), porte (\texttt{Door}), fenêtre (\texttt{Window}).
    \item \textbf{Tête Sémantique (12 classes) :} Fond, Living Room, Bedroom, Kitchen, Bathroom, Entry, Storage, Balcony, Corridor, Utility, Office, Undefined.
\end{itemize}

La perte d'entraînement conjointe associe la Cross-Entropy pondérée (\emph{Weighted Cross-Entropy}) pour mitiger le déséquilibre des parois minces, et la perte Dice multimodale :
\begin{equation}
    \mathcal{L}_{\text{struct}} = 0.5 \cdot \mathcal{L}_{\text{WCE}} + 0.5 \cdot \mathcal{L}_{\text{Dice}} \quad ; \quad \mathcal{L}_{\text{room}} = 0.5 \cdot \mathcal{L}_{\text{WCE}} + 0.5 \cdot \mathcal{L}_{\text{Dice}}
\end{equation}
\begin{equation}
    \mathcal{L}_{\text{multi-task}} = \mathcal{L}_{\text{struct}} + 0.6 \cdot \mathcal{L}_{\text{room}}
\end{equation}

\newpage

% =========================================================================
% CHAPITRE 3: PROTOCOLE EXPÉRIMENTAL & DONNÉES
% =========================================================================
\chapter{Protocole Expérimental \& Données}

\section{Jeux de Données et Découpage Standardisé}
\begin{enumerate}
    \item \textbf{CMP Facades Dataset :} Composé de 606 paires d'images de façades architecturales et de leurs étiquetages sémantiques. Le découpage officiel comprend 400 paires en apprentissage, 100 en validation et 106 en test indépendant.
    \item \textbf{CubiCasa5K :} 5\,000 plans d'étage d'architectes annotés au format vectoriel SVG. Le partitionnement officiel adopté respecte strictement la séparation standardisée : 4\,200 plans d'entraînement (\texttt{train}), 400 plans de validation (\texttt{val}), et 400 plans de test (\texttt{test}).
\end{enumerate}

\section{Audit et Résolution du Bug de Rasterisation SVG}
L'audit du pipeline de données initial a permis d'isoler une anomalie critique : l'utilisation d'une recherche de sous-chaîne naïve \texttt{"Wall" in class\_name} entraînait l'assimilation erronée des éléments de mobilier \texttt{"WallCabinet"} comme des murs porteurs, et la traversée récursive de l'arbre SVG écrasait les portes et fenêtres imbriquées au profit des contours de murs.
La correction implémentée applique un filtrage strict par nom exact et sépare les cloisons externes des cloisons internes via l'analyse de contiguïté avec les espaces extérieurs.

\section{Reproductibilité et Repères d'Entraînement}
Toutes les expériences sont exécutées sous contrainte déterministe avec verrouillage de la graine pseudo-aléatoire à $42$ :
\begin{lstlisting}[language=Python]
def seed_everything(seed=42):
    random.seed(seed)
    np.random.seed(seed)
    torch.manual_seed(seed)
    torch.cuda.manual_seed_all(seed)
    torch.backends.cudnn.deterministic = True
\end{lstlisting}

\newpage

% =========================================================================
% CHAPITRE 4: RÉSULTATS EXPÉRIMENTAUX & ANALYSE CRITIQUE
% =========================================================================
\chapter{Résultats Expérimentaux \& Analyse Critique}

\section{Évaluation de la Synthèse Générative de Façades}
Le tableau~\ref{tab:facades_benchmark} synthétise les performances obtenues sur les 106 images du jeu de test CMP Facades.

\begin{table}[H]
\centering
\caption{Benchmark Comparatif des Modèles Génératifs (CMP Facades Test Split --- $N = 106$)}
\label{tab:facades_benchmark}
\begin{tabular}{lcccc}
\toprule
\textbf{Architecture} & \textbf{FID $\downarrow$} & \textbf{KID $(\times 10^3) \downarrow$} & \textbf{SSIM $\uparrow$} & \textbf{$\mathcal{L}_1 \downarrow$} \\
\midrule
Baseline (100 ep, TransposeConv, $\mathcal{L}_1=100$) & \textbf{351.70} & 149.89 $\pm$ 10.27 & 0.1226 & 0.3986 \\
Enhanced (150 ep, Nearest+Conv, $\mathcal{L}_1=100$)  & 371.59 & 103.79 $\pm$ 8.89  & \textbf{0.1430} & \textbf{0.3918} \\
Perceptual (100 ep, LPIPS+SN, $\mathcal{L}_1=30$)     & 369.94 & \textbf{60.40 $\pm$ 5.52} & 0.1058 & 0.3968 \\
\bottomrule
\end{tabular}
\end{table}

\subsection{Discussion Scientifique des Désaccords entre Métriques}
Une analyse rigoureuse impose de reconnaître que les métriques objectives divergent :
\begin{enumerate}
    \item \textbf{Métrique FID \& Biais d'Échantillonnage :} Le score FID absolu (350--370) s'avère élevé en raison de la dimensionnalité de l'espace Inception ($d = 2048$) vis-à-vis du faible effectif d'échantillons ($N = 106$), rendant l'estimation de la matrice de covariance de Fréchet sous-dimensionnée. Le KID (Kernel Inception Distance), qui fournit un estimateur non biaisé du $\text{MMD}^2$, classe quant à lui le modèle perceptuel en tête ($60.40 \times 10^{-3}$).
    \item \textbf{Pénalisation de la Netteté par la Perte $\mathcal{L}_1$ et le SSIM :} Le modèle \emph{Enhanced} obtient les meilleurs scores SSIM ($0.1430$) et $\mathcal{L}_1$ ($0.3918$) parce que la minimisation exclusive de l'erreur absolue force le modèle à converger vers la moyenne conditionnelle pixel-par-pixel, ce qui produit un flou moyen statistiquement récompensé mais visuellement pauvre. Le modèle perceptuel, entraîné avec une composante LPIPS, affiche un SSIM ($0.1058$) et un $\mathcal{L}_1$ ($0.3968$) inférieurs, mais une distance perceptuelle et un KID nettement plus favorables.
    \item \textbf{Justification du Choix de Production :} Le modèle perceptuel a été sélectionné pour le moteur d'inférence en raison de son acuité visuelle : suppression visible du motif d'échiquier grâce aux blocs Nearest+Conv, stabilisation du discriminateur par normalisation spectrale, encadrements de fenêtres nets et restitution des textures de pierre architecturale.
\end{enumerate}

\section{Évaluation de la Segmentation de Plans d'Étage}
Le tableau~\ref{tab:floorplan_benchmark} expose les résultats empiriques mesurés sur l'intégralité des 400 plans du jeu de test CubiCasa5K non vus.

\begin{table}[H]
\centering
\caption{Benchmark sur le Jeu de Test CubiCasa5K ($400$ plans complets non vus)}
\label{tab:floorplan_benchmark}
\begin{tabular}{lcc}
\toprule
\textbf{Métrique} & \textbf{Baseline ResNet-34} & \textbf{MultiTask UNet++ (Aligné)} \\
\midrule
IoU Murs Fusionnés (Merged Wall IoU) & 29.08\% & \textbf{74.32\%} \\
IoU Murs Extérieurs & N/A & \textbf{73.89\%} \\
IoU Murs Intérieurs & N/A & \textbf{62.34\%} \\
IoU Portes (Door IoU) & 7.91\% & \textbf{62.38\%} \\
IoU Fenêtres (Window IoU) & 7.28\% & \textbf{73.28\%} \\
mIoU Structurelle Premier Plan & 14.76\% & \textbf{69.99\%} \\
mIoU Sémantique des Pièces (11 classes) & N/A & \textbf{54.72\%} \\
Exactitude Pixel des Pièces & N/A & \textbf{85.85\%} (Réf. papier: 82.7\%) \\
Rappel Portes à 20px (Object Recall) & 53.10\% & \textbf{96.75\%} \\
Rappel Fenêtres à 20px (Object Recall) & 49.14\% & \textbf{90.59\%} \\
IoU Empreinte Habitable (Living Area) & N/A & \textbf{74.04\%} \\
Erreur Absolue Moyenne Nb Pièces (MAE) & N/A & \textbf{2.61} pièces \\
Exactitude Classification Murs Extérieurs & 42.88\% & \textbf{91.07\%} (Préc: 94.2\%, Rap: 92.4\%) \\
Erreur Relative d'Échelle Physique & 44.15\% & \textbf{10.14\%} \\
Taux de Déclenchement Auto-Échelle & 100.0\% & 100.0\% \\
Erreur Relative de Périmètre & 81.82\% & \textbf{45.08\%} \\
Taux de Coque 3D Étanche (Watertight) & 9.75\% & \textbf{96.25\%} \\
\bottomrule
\end{tabular}
\end{table}

\subsection{Diagnostic et Résolution du Problème Fondamental d'Alignement Données-Masques}
L'historique de développement démontre une avancée méthodologique cruciale :
\begin{itemize}
    \item \textbf{Cause Racine Identifiée (Désynchronisation Vectorielle-Raster) :} L'analyse fine a révélé un décalage géométrique systématique dans le pipeline d'ingestion initial : les coordonnées des polygones vectoriels SVG bruts de CubiCasa5K n'étaient pas correctement recalées sur le canevas matriciel $512 \times 512$ letterboxé. Ce décalage induisait une pénalisation artificielle où les cloisons réelles et prédites ne se superposaient jamais exactement, bridant l'IoU des murs à $22.44\%$ et des portes à $10.71\%$.
    \item \textbf{Résolution et Preuve d'Apprentissage :} Après réécriture rigoureuse de la mise à l'échelle géométrique dans \texttt{src/floorplan/dataset.py} et mise en cache \texttt{.npz}, un test de sur-apprentissage contrôlé a immédiatement confirmé la convergence exacte ($>90\%$ d'IoU). Le ré-entraînement complet sur 25 époques a validé ce bond qualitatif sur les 400 plans de test : $+51.88$ pts sur les murs ($74.32\%$), $+51.67$ pts sur les portes ($62.38\%$), et $+49.90$ pts sur les fenêtres ($73.28\%$).
    \item \textbf{Comparaison aux Résultats Publiés :} L'exactitude pixel des pièces atteint $85.85\%$, surpassant le benchmark officiel de l'article CubiCasa5K ($82.7\%$). Le mIoU sémantique des pièces s'établit à $54.72\%$, proche du score publié ($57.5\%$), tout en opérant sans aucun post-traitement de nettoyage manuel.
\end{itemize}

\newpage

% =========================================================================
% CHAPITRE 5: MLOPS & ARCHITECTURE LOGICIELLE
% =========================================================================
\chapter{Cycle de Vie MLOps \& Architecture Logicielle}

\section{Versioning des Données et Modèles (DVC)}
Afin de respecter scrupuleusement la règle d'or interdisant le stockage de fichiers binaires volumineux sur Git, l'intégralité des données brutes et des checkpoints est versionnée via Data Version Control (DVC). Seuls les pointeurs textuels légers (\texttt{.dvc}) et les règles \texttt{.gitignore} sont intégrés aux commits Git.

\section{Traçabilité et Registre des Modèles (MLflow)}
Le serveur MLflow assure le suivi automatique de chaque itération :
\begin{itemize}
    \item Enregistrement des scalaires à chaque époque : pertes d'entraînement, mIoU, IoU par classe, taux d'apprentissage (\texttt{lr}).
    \item Sauvegarde automatique d'artefacts visuels : une planche comparative de 6 plans fixes de validation est générée et loggée toutes les 5 époques.
\end{itemize}

\section{Moteur d'Inférence ONNX Runtime \& Parité Numérique}
Pour découpler le service en production des dépendances lourdes de PyTorch et optimiser le débit d'exécution, le modèle génératif a été converti au format ONNX Runtime. Sur station GPU locale (NVIDIA GeForce RTX 3050 Laptop, 4 Go VRAM, CPU AMD Ryzen 5), la latence de l'inférence tensorielle UNet++ s'établit à $45\text{ ms}$, tandis que le pipeline complet de bout en bout (segmentation multi-tâches, vectorisation polygonale Shapely, synthèse de texture de façade cGAN et extrusion 3D OBJ/MTL) s'exécute en $2.03\text{ s}$ ($2\,031.9\text{ ms}$) en moyenne. La parité numérique entre PyTorch et ONNX Runtime est rigoureusement confirmée :
\begin{equation}
    \Delta_{\max} = \max_{x \in \mathcal{D}} | \mathcal{M}_{\text{PyTorch}}(x) - \mathcal{M}_{\text{ONNXRuntime}}(x) | = 6.31 \times 10^{-6}
\end{equation}
Cette divergence négligeable garantit une stricte équivalence numérique en production.

\section{Déploiement Conteneurisé Docker Compose}
L'architecture est orchestrée par \texttt{docker-compose.yml} en deux micro-services découplés :
\begin{itemize}
    \item \textbf{Backend FastAPI (Port 8000) :} Expose les endpoints \texttt{/health}, \texttt{/metrics}, \texttt{/predict/floorplan} et \texttt{/predict/render}. Intègre un middleware mesurant la latence réelle retournée dans l'en-tête HTTP \texttt{X-Inference-Latency-ms}.
    \item \textbf{Frontend Nginx (Port 3000) :} Sert le bundle de production React 18 / Three.js compilé par Vite, et intègre un proxy inverse redirigeant les requêtes d'API vers \texttt{http://backend:8000}.
\end{itemize}

\newpage

% =========================================================================
% CHAPITRE 6: GÉNÉRALISATION & ANALYSE DES MODES D'ÉCHEC
% =========================================================================
\chapter{Généralisation \& Analyse des Modes d'Échec}

\section{Tests de Robustesse Synthétiques (10 Plans de Stress Hors-Domaine)}
Afin de sonder les limites de robustesse et la sensibilité aux dérives de distribution (\emph{domain shifts}), 10 plans de stress synthétiques couvrant des conventions graphiques variées ont été soumis au pipeline. Le tableau~\ref{tab:external_benchmark} confronte les prédictions au nombre exact de pièces tracées en vérité terrain.

\begin{table}[H]
\centering
\caption{Synthèse des Tests de Robustesse Synthétiques ($10$ cas hors CubiCasa5K)}
\label{tab:external_benchmark}
\begin{tabular}{clccccl}
\toprule
\textbf{\#} & \textbf{Cas de Stress} & \textbf{GT} & \textbf{Préd.} & \textbf{$\Delta R$} & \textbf{Coque} & \textbf{Statut Rigoureux / Mode de Défaillance} \\
\midrule
1 & CAD N\&B épuré & 4 & 0 & -4 & Échec & ÉCHEC : Aucune pièce détectée (baies ouvertes) \\
2 & Brochure pastel & 3 & 0 & -3 & Échec & ÉCHEC : Dérive spectrale des textures de sol \\
3 & Esquisse main levée & 3 & 1 & -2 & Valide & SOUS-SEGMENTÉ : 1 pièce (flou de traits ondulés) \\
4 & Plan encombré mobilier & 4 & 2 & -2 & Valide & SOUS-SEGMENTÉ : 2 pièces (occlusion par mobilier) \\
5 & Cotes \& textes techniques & 4 & 0 & -4 & Échec & ÉCHEC : Textes confondus avec cloisons \\
6 & Basse résolution (140px) & 4 & 7 & +3 & Incomplète & SUR-SEGMENTÉ : Artefacts de compression flous \\
7 & Loft ouvert moderne & 2 & 0 & -2 & Échec & ÉCHEC : Absence de cloisons physiques fermées \\
8 & Logement allongé (couloir) & 5 & 0 & -5 & Échec & ÉCHEC : Cloisonnement longitudinal non refermé \\
9 & CAD sombre inversé & 4 & 10 & +6 & Incomplète & DÉFAILLANCE : Inversion contraste fond/murs \\
10 & Murs en diagonale (45\textdegree) & 2 & 0 & -2 & Échec & ÉCHEC : Snapping rectilinéaire orthogonal \\
\bottomrule
\end{tabular}
\end{table}

\section{Analyse Approfondie des Dérives de Domaine}
Cette campagne de stress tests dissipe toute ambiguïté quant à la capacité de généralisation :
\begin{enumerate}
    \item \textbf{Échecs de Détection Complète (6/10 cas) :} Pour 6 des 10 plans, aucune pièce n'est détectée. Les versions antérieures qualifiaient à tort ces sorties vides de << coques valides >>. Le système renvoie désormais un statut explicite \texttt{NO\_ROOMS\_DETECTED}, signalant à l'API et à l'utilisateur que l'enveloppe ne contient aucun espace habitable.
    \item \textbf{Sensibilité Extrême aux Ruptures de Style :} L'entraînement exclusif sur CubiCasa5K (styles scandinaves vectorisés standardisés) rend le réseau vulnérable aux styles de dessin alternatifs : les fonds noirs inversés activent faussement le décodeur convolutif ($+6$ pièces parasites), et le flou de sous-échantillonnage fragmente les parois minces en cellules orphelines.
    \item \textbf{Nécessité de Garde-Fous Interactifs :} Ces constats justifient pleinement l'intégration dans l'application web de l'outil de calibration d'échelle à 2 points et d'une porte de plausibilité architecturale, avertissant l'utilisateur lorsque l'agencement extrait s'écarte des normes physiques résidentielles.
\end{enumerate}

\newpage

% =========================================================================
% CONCLUSION & PERSPECTIVES
% =========================================================================
\chapter{Conclusion \& Perspectives}

Le projet \textbf{ArchiGen 3D} concrétise l'industrialisation réussie d'une solution complète d'IA architecturale, reliant vision par ordinateur, géométrie computationnelle, génération d'images antagoniste et développement full-stack moderne. 

L'entraînement rigoureux du modèle MultiTask UNet++ pendant 60 époques sur CubiCasa5K a permis d'atteindre un taux d'étanchéité de coque de $99.75\%$, un IoU d'empreinte habitable de $61.91\%$ et une exactitude de classification des pièces de $71.11\%$. Côté génération, l'architecture Pix2Pix intégrant normalisation spectrale et perte perceptuelle LPIPS produit des textures de façades d'une grande finesse architecturale sous ONNX Runtime.

\textbf{Perspectives d'évolution :}
\begin{enumerate}
    \item Remplacement de la segmentation purement matricielle des cloisons par un graphe relationnel de pièces (\emph{Room Layout Graphs} avec Graph Neural Networks - GNN).
    \item Intégration de modèles de diffusion latente conditionnés par la structure géométrique (ControlNet / Stable Diffusion) pour offrir une variété de styles de rendu encore plus étendue.
\end{enumerate}

\end{document}
